\section{Introduction}~\label{intro}
% Robots operating in real-world environments often have incomplete or unreliable information about the current state due to perception errors~\citep{kaipa2016addressing}, execution failures~\citep{kaelbling1998planning}, and unexpected environmental changes~\citep{kurniawati2016online}. If such uncertainty is not resolved, the robot may make incorrect decisions and fail the task. In practice, relying on an external agent to recognize the need for assistance requires continuous monitoring of robot execution. The robot must therefore be able to detect uncertainty and initiate queries to obtain additional information before proceeding~\citep{dragan2013legibility, goodrich2008human}. However, frequent queries can incur substantial costs in time and workload~\citep{chen2014human}. The robot must therefore determine \textit{when} to request external feedback and 
Robots often encounter uncertainty during task planning due to perception errors~\citep{kaipa2016addressing} and action failures~\citep{kaelbling1998planning}. For example, a robot may be uncertain whether a tomato is ripe or whether a pick action succeeded. The robot can proactively ask an external agent to provide missing information or select the next action to reduce the need for continuous monitoring by the external agent. However, too few queries can leave uncertainty unresolved and lead to task failure, whereas frequent queries increase interaction costs, such as response latency, token use for a vision language model (VLM), and human workload. A further challenge is that simply increasing the number of queries does not necessarily improve task success. The robot may still fail if queries come late or do not resolve the uncertainty. The robot must therefore determine \textit{when} to ask for help and \textit{what} information to request for task success while limiting the number of queries.

% Recent uncertainty-aware systems selectively request external feedback during task execution. Existing approaches commonly trigger queries based on prediction confidence, including conformal prediction~\citep{ren2023robots, liang2024introspective, mullen2025lbap}, or on the expected effect of a query on task reward~\citep{burks2021collaborative, anand2025adaptive}. However, a confident action prediction does not ensure that the underlying state estimate is correct. Under reward-based criteria, the robot may also execute an action without querying even when state uncertainty remains. Both approaches can therefore miss opportunities to verify the state information needed for successful execution, leading to task failure. We use state uncertainty to determine when verification is needed and expected uncertainty reduction to select what to ask.
% Existing system-initiative querying methods often use action prediction confidence or expected return to guide querying decisions. For example, KnowNo uses conformal prediction sets to determine when to request an action choice~\citep{ren2023robots}, while IntroPlan combines conformal prediction with retrieved reasoning examples for candidate evaluation~\citep{liang2024introspective}. Yet both methods may proceed without querying if incorrect perception produces a confident action prediction. Even when assistance is requested, the candidate set may not include a correct action. POMDP-based approaches~\citep{burks2021collaborative,burks2023harps} model candidate queries as planning actions and determine when and what to query by maximizing expected cumulative reward. However, maximizing cumulative rewards for task performance does not necessarily yield efficient querying, because asking fewer questions can increase the risk of failure, while asking more does not always improve task success.
Existing methods often use confidence in predicted actions or expected reward return to decide when and what to query. Methods based on conformal prediction (CP) query an external agent if multiple candidate actions remain~\citep{ren2023robots,liang2024introspective}. However, a single candidate does not ensure a correct state estimate, and the prediction set may exclude the correct action. Approaches based on a partially observable Markov decision process (POMDP) model queries as planning actions~\citep{burks2021collaborative,burks2023harps} and select them according to expected cumulative reward. However, neither approach explicitly determines whether the current state is sufficiently certain. Consequently, the robot may plan with insufficient state information or ask unnecessary queries.


% In this paper, we present a belief-based method for deciding when to request external feedback and what to ask. We formulate task planning as a POMDP~\citep{kaelbling1998planning} and maintain a knowledge base containing symbolic facts and STRIPS-style action models~\citep{davis1993knowledge,fikes1971strips}. After each action, the robot uses these facts and action models to generate reachable states and constructs a weighted belief using the resulting observation. If belief confidence falls below a threshold, the robot asks an external agent to verify the state fact with the greatest expected entropy reduction. The robot updates the belief after each response and continues querying until the confidence reaches the threshold. The robot then updates the facts in the knowledge base and resumes planning and execution.
% In this paper, we present a method that uses uncertainty about the current symbolic state to determine when to query and what to ask. We formulate task planning as a POMDP~\citep{kaelbling1998planning} and use symbolic facts and action models~\citep{davis1993knowledge,fikes1971strips} to construct a weighted belief over reachable states. To determine when to query, the robot computes belief confidence from normalized entropy, which decreases as probability concentrates on a single state. The robot queries an external agent if belief confidence falls below a threshold. To determine what to query, the robot selects the state fact with the greatest expected reduction in belief entropy to reduce the number of queries needed to reach the confidence threshold. Each response provides evidence for updating the belief. Once confidence reaches the threshold, the robot updates the knowledge base with the maximum a posteriori state estimate and resumes planning and execution.

In this paper, we present a method that uses uncertainty about the current symbolic state to determine when to query and what to ask. We formulate task planning as a POMDP~\citep{kaelbling1998planning} and use symbolic facts and action models~\citep{davis1993knowledge,fikes1971strips} to construct a weighted belief over reachable states. To determine when to query, the robot computes belief confidence from normalized entropy, which decreases as probability concentrates on a single state. The robot queries an external agent if the confidence falls below a threshold. To determine what to query, the robot selects the fact with the greatest expected reduction in belief entropy. Each response rules out states that conflict with the answer. We use the maximum a posteriori (MAP) state to update the knowledge base and resume planning. We construct the next frontier from the MAP state and the next action.



% We evaluate the method in two robotic domains. In simulation with correct oracle responses, we compare our method with conformal prediction based and reward-based strategies and analyze the effects of query timing and query selection. We then conduct physical robot experiments with feedback from a vision-language model (VLM) and human participants, comparing our method with the conformal prediction based method. Interviews with the five participants provide qualitative insights into their experiences with both methods. Our method achieves higher task success rates than all evaluated baselines in simulation and than the selected baseline with VLM and human feedback on physical robots. Ablation results indicate that query timing primarily affects task success, while query selection reduces the number of queries required.

% We evaluate our method on two robotic tasks in partially observable environments with sensing errors. In simulations with an oracle agent, we compare our method against two query strategies based on conformal prediction and expected task reward, respectively. We then conduct physical robot experiments using a VLM and human participants as external agents to compare our method with a baseline based on conformal prediction. Our method achieves the highest success rates in both simulation and physical robot experiments with VLM and human feedback. Ablation experiments show that query timing mainly affects task success, while our query selection method reduces the number of queries with the same timing criterion. We also vary the confidence threshold to examine the trade-off between task success and query count. Finally, interviews with all five participants help explain the differences in task performance between the two methods.
We evaluate our method on two robotic tasks in partially observable environments with sensing errors. In simulations with an oracle agent, we first compare our system against existing proactive querying systems based on CP and expected task reward. We then conduct physical robot experiments with VLM and human feedback to compare our method with a CP baseline. Our method achieves the highest success rates among the evaluated systems across all baseline experiments. We then conduct ablation experiments with the same POMCP planner to distinguish the effects of query timing and selection on task success and query count. Ablation experiments show that query timing based on confidence maintains high success across selection strategies, while query selection based on expected entropy reduction reduces the number of queries under the same timing criterion. We also vary the confidence threshold to examine the tradeoff between task success and query count. Finally, interviews with all five participants help explain the differences in task performance between the two methods.

% The contributions of this work are as follows:
% \begin{itemize}
% \item We determine when to query based on uncertainty in the current belief and show that query timing has a substantial effect on task success.
% \item We choose what to query using expected entropy reduction and show that query selection reduces query counts compared to the baselines.
% \item We evaluate the proposed method in simulation and on physical robots with VLM and human feedback, demonstrating higher task success than the evaluated baselines.
% \end{itemize}
Our contributions are as follows:
\begin{itemize}
\item We determine when to query from belief uncertainty to improve task success.
\item We select what to query by expected entropy reduction, reducing query counts under a fixed timing criterion.
\item We demonstrate higher task success than the evaluated baselines in simulation and on physical robots with VLM and human feedback.
\end{itemize}